\documentclass[11pt,a4paper]{article}
\usepackage[utf8]{inputenc}
\usepackage[T1]{fontenc}
\usepackage{amsmath,amssymb,amsthm}
\usepackage{listings}
\usepackage{geometry}
\usepackage{hyperref}
\geometry{margin=2.5cm}

\newtheorem{theorem}{Theorem}
\newtheorem{lemma}[theorem]{Lemma}

\lstset{
  language=C++,
  basicstyle=\ttfamily\small,
  keywordstyle=\bfseries,
  breaklines=true,
  numbers=left,
  numberstyle=\tiny,
  frame=single,
  tabsize=2
}

\title{IOI 2023 --- Longest Trip}
\author{}
\date{}

\begin{document}
\maketitle

\section{Problem Statement}

$N$ landmarks are connected by roads.  The graph satisfies a
\emph{density} condition: for any triple of nodes, at least $D$ of
the three possible edges exist ($D \in \{1, 2, 3\}$).  Using an
interactive function $\texttt{are\_connected}(A, B)$ (which tests
whether any edge exists between sets $A$ and $B$), find the longest
simple path.

\section{Solution}

\subsection{Case $D = 3$: complete graph}

Every pair is connected, so any permutation of all $N$ nodes is a
Hamiltonian path.

\subsection{Case $D = 2$}

Any three nodes span at least 2 edges.  A Hamiltonian path always
exists.

\subsection{Case $D = 1$: general}

\begin{lemma}\label{lem:merge}
When adding node $i$ that is not adjacent to any endpoint of path $P$
or path $Q$, the density condition forces $P$'s last node and $Q$'s
last node to be adjacent.  Hence we can concatenate $P$ and $Q$.
\end{lemma}

\begin{proof}
Consider the triple ($i$, $P.\mathrm{back}$, $Q.\mathrm{back}$).  We
tested that $i$--$P.\mathrm{back}$ and $i$--$Q.\mathrm{back}$ are
non-edges.  Since at least 1 edge must exist among the three, the edge
$P.\mathrm{back}$--$Q.\mathrm{back}$ exists.
\end{proof}

\paragraph{Algorithm.}
Maintain two paths $P$ and $Q$ (as deques).  For each new node $i$:
\begin{enumerate}
  \item If $i$ connects to an endpoint of $P$, extend $P$.
  \item Else if $i$ connects to an endpoint of $Q$, extend $Q$.
  \item Else by Lemma~\ref{lem:merge}, merge $P$ and $Q$ (appending
        $Q$ reversed to $P$), then start a new $Q = \{i\}$.
\end{enumerate}

After all nodes are processed, try to merge $P$ and $Q$ by testing
all four endpoint combinations.  If no endpoint connection exists,
use binary search to find an internal connection point (costing
$O(\log N)$ additional queries) and splice the paths together.

\section{Implementation}

\begin{lstlisting}
#include <bits/stdc++.h>
using namespace std;

bool are_connected(vector<int> A, vector<int> B);

vector<int> longest_trip(int N, int D) {
    if (N == 1) return {0};

    deque<int> P, Q;
    P.push_back(0);
    if (N > 1) Q.push_back(1);

    for (int i = 2; i < N; i++) {
        if (are_connected({i}, {(int)P.back()})) {
            P.push_back(i);
        } else if (are_connected({i}, {(int)P.front()})) {
            P.push_front(i);
        } else if (!Q.empty() && are_connected({i}, {(int)Q.back()})) {
            Q.push_back(i);
        } else if (!Q.empty() && are_connected({i}, {(int)Q.front()})) {
            Q.push_front(i);
        } else {
            // By density D >= 1: P.back()--Q.back() must be an edge
            while (!Q.empty()) {
                P.push_back(Q.back());
                Q.pop_back();
            }
            Q.push_back(i);
        }
    }

    if (Q.empty())
        return vector<int>(P.begin(), P.end());

    // Try all 4 endpoint pairs
    if (are_connected({(int)P.back()}, {(int)Q.front()})) {
        for (auto x : Q) P.push_back(x);
        return vector<int>(P.begin(), P.end());
    }
    if (are_connected({(int)P.back()}, {(int)Q.back()})) {
        while (!Q.empty()) { P.push_back(Q.back()); Q.pop_back(); }
        return vector<int>(P.begin(), P.end());
    }
    if (are_connected({(int)P.front()}, {(int)Q.front()})) {
        vector<int> res;
        while (!P.empty()) { res.push_back(P.back()); P.pop_back(); }
        for (auto x : Q) res.push_back(x);
        return res;
    }
    if (are_connected({(int)P.front()}, {(int)Q.back()})) {
        while (!Q.empty()) { P.push_front(Q.back()); Q.pop_back(); }
        return vector<int>(P.begin(), P.end());
    }

    // No endpoint connection: binary search for internal link
    vector<int> pv(P.begin(), P.end()), qv(Q.begin(), Q.end());
    if (!are_connected(pv, qv))
        return (pv.size() >= qv.size()) ? pv : qv;

    // Binary search on P for a node connected to Q
    int lo = 0, hi = (int)pv.size() - 1, p_idx = 0;
    while (lo < hi) {
        int mid = (lo + hi) / 2;
        vector<int> sub(pv.begin() + lo, pv.begin() + mid + 1);
        if (are_connected(sub, qv))
            hi = mid;
        else
            lo = mid + 1;
    }
    p_idx = lo;

    // Binary search on Q for the partner
    lo = 0; hi = (int)qv.size() - 1;
    int q_idx = 0;
    while (lo < hi) {
        int mid = (lo + hi) / 2;
        vector<int> sub(qv.begin() + lo, qv.begin() + mid + 1);
        if (are_connected({pv[p_idx]}, sub))
            hi = mid;
        else
            lo = mid + 1;
    }
    q_idx = lo;

    // Splice: ...P[p_idx] -- Q[q_idx]...
    vector<int> result;
    for (int i = p_idx + 1; i < (int)pv.size(); i++) result.push_back(pv[i]);
    for (int i = 0; i <= p_idx; i++) result.push_back(pv[i]);
    for (int i = q_idx; i < (int)qv.size(); i++) result.push_back(qv[i]);
    for (int i = 0; i < q_idx; i++) result.push_back(qv[i]);
    return result;
}
\end{lstlisting}

\section{Complexity}

\begin{itemize}
  \item \textbf{Queries during construction:} $O(N)$ --- at most 4 queries
        per node for endpoint checks.
  \item \textbf{Queries for merging:} $O(\log N)$ for binary search on each
        path.  Total: $O(\log^2 N)$.
  \item \textbf{Overall:} $O(N + \log^2 N) = O(N)$ queries.
  \item \textbf{Space:} $O(N)$.
\end{itemize}

\end{document}
